\documentclass{article}

\usepackage[british]{babel}
\usepackage[T1]{fontenc}
\usepackage{lmodern}
\usepackage{microtype}

\usepackage{amsmath}
\usepackage{amssymb}
\usepackage{amsthm}

\usepackage{fullpage}
\usepackage{parskip}
\usepackage{fancyhdr}
\usepackage{enumitem}

\usepackage{xcolor}
\usepackage{hyperref}
\usepackage[capitalise,nameinlink]{cleveref}

\hypersetup{
  colorlinks=true,
  allcolors=blue!50!black,
  pdftitle={Odd degree extensions},
  pdfauthor={Joe Lamond and Joshua Munchau},
}

\pagestyle{fancy}
\fancyhf{}
\fancyhead[L]{Odd degree extensions}
\fancyhead[R]{Joe Lamond and Joshua Munchau}
\fancyfoot[C]{\thepage}
\setlength{\headheight}{14pt}
\setlength{\headsep}{10pt}

\newtheorem{theorem}{Theorem}
\newtheorem*{remark}{Remark}

\DeclareMathOperator{\Gal}{Gal}
\DeclareMathOperator{\Aut}{Aut}

\begin{document}

\section*{Odd degree extensions}
This is Step~1 of the proof sketch in the \emph{Outline of formalisation}.
\begin{theorem}\label{thm:odd-degree-extensions}
Suppose $R$ is a real closed field, and $K/R$ is a finite extension of
odd degree. Then $[K:R]=1$.
\end{theorem}
(For simplicity, we shall regard $R$ as a subfield of $K$, though in the Lean
file the extension is treated as a ring homomorphism $R\to K$, and the goal is
to show that it is bijective.)
\begin{proof}
Let $\alpha\in K$. Since $K$ is finite over $R$, the element $\alpha$
is algebraic over $R$. Let $p\in R[t]$ be its minimal polynomial. By the tower
law, $[R(\alpha):R]$ divides $[K:R]$. Since $[K:R]$ is odd, $[R(\alpha):R]$ is
odd. Thus $p$ has odd degree.

By the definition of real closed field, $p$ has a root, say $c\in R$. Thus $t-c$
divides $p$, and since $p$ is monic and irreducible, we have $p=t-c$. Now 
$p(\alpha)=0$, so that $\alpha=c\in R$. Since $\alpha$ was arbitrary,
we conclude that $K=R$, hence $[K:R]=1$.
\end{proof}

\end{document}
